
\documentclass{article}
\usepackage{colortbl}
\usepackage{makecell}
\usepackage{multirow}
\usepackage{supertabular}

\begin{document}

\newcounter{utterance}

\centering \large Interaction Transcript for game `hot\_air\_balloon', experiment `air\_balloon\_survival\_en\_negotiation\_hard', episode 3 with deepseek{-}chat{-}v3.1{-}t1.0.
\vspace{24pt}

{ \footnotesize  \setcounter{utterance}{1}
\setlength{\tabcolsep}{0pt}
\begin{supertabular}{c@{$\;$}|p{.15\linewidth}@{}p{.15\linewidth}p{.15\linewidth}p{.15\linewidth}p{.15\linewidth}p{.15\linewidth}}
    \# & $\;$A & \multicolumn{4}{c}{Game Master} & $\;\:$B\\
    \hline

    \theutterance \stepcounter{utterance}  
    & & \multicolumn{4}{p{0.6\linewidth}}{
        \cellcolor[rgb]{0.9,0.9,0.9}{
            \makecell[{{p{\linewidth}}}]{
                \texttt{\tiny{[P1$\langle$GM]}}
                \texttt{You are participating in a collaborative negotiation game.} \\
\\ 
\texttt{Together with another participant, you must agree on a single set of items that will be kept. Each of you has your own view of how much each item matters to you (importance). You do not know how the other participant values the items. Additionally, you are given the effort each item demands.} \\
\texttt{You may only agree on a set if the total effort of the selected items does not exceed a shared limit:} \\
\\ 
\texttt{LIMIT = 3681} \\
\\ 
\texttt{Here are the individual item effort values:} \\
\\ 
\texttt{Item effort = \{"C17": 700, "B41": 586, "A71": 505, "C00": 116, "B35": 664, "A77": 940, "A03": 387, "A04": 392, "C54": 209, "C01": 571, "A87": 4, "B32": 285, "B86": 651, "B33": 613, "A10": 740\}} \\
\\ 
\texttt{Here is your personal view on the importance of each item:} \\
\\ 
\texttt{Item importance values = \{"C17": 138, "B41": 583, "A71": 868, "C00": 822, "B35": 783, "A77": 65, "A03": 262, "A04": 121, "C54": 508, "C01": 780, "A87": 461, "B32": 484, "B86": 668, "B33": 389, "A10": 808\}} \\
\\ 
\texttt{Goal:} \\
\\ 
\texttt{Your goal is to negotiate a shared set of items that benefits you as much as possible (i.e., maximizes total importance to YOU), while staying within the LIMIT. You are not required to make a PROPOSAL in every message {-} you can simply negotiate as well. All tactics are allowed!} \\
\\ 
\texttt{Interaction Protocol:} \\
\\ 
\texttt{You may only use the following structured formats in a message:} \\
\\ 
\texttt{PROPOSAL: \{'A', 'B', 'C', …\}} \\
\texttt{Propose keeping exactly those items.} \\
\texttt{REFUSE: \{'A', 'B', 'C', …\}} \\
\texttt{Explicitly reject opponent's proposal.} \\
\texttt{ARGUMENT: \{'...'\}} \\
\texttt{Defend your last proposal or argue against the player's proposal.} \\
\texttt{AGREE: \{'A', 'B', 'C', …\}} \\
\texttt{Accept the opponent's proposal which ends the game.} \\
\texttt{STRATEGIC REASONING: \{'...'\}} \\
\texttt{	Describe your strategic reasoning or anticipation explaining your choice of action. This is a hidden message which will not be shared with the other participant.} \\
\\ 
\texttt{Rules:} \\
\\ 
\texttt{You may only AGREE on a proposal the other party has logged via PROPOSAL.} \\
\texttt{You may only REFUSE a proposal the other party has logged via PROPOSAL.} \\
\texttt{Total effort of any PROPOSAL or AGREE set must be ≤ LIMIT.} \\
\texttt{Do NOT reveal your hidden importance scores.} \\
\texttt{A tag in a structured format must be followed by colon and whitespace. The argument must be a python set containing 0 or more strings.} \\
\texttt{So, it must be of the form TAG: \{...\}} \\
\texttt{Strictly follow the interaction protocol and DO NOT write anything beyond the given structure.} \\
\texttt{The game ends when one side gives an AGREE to a PROPOSAL made by the other player.} \\
\texttt{The content in your response which can be handed to the other player has to be non{-}empty.} \\
\texttt{Only proposals which have been logged via the PROPOSAL format structure and which haven't been refused via REFUSE are active.} \\
\texttt{You must include the ARGUMENT format at least once somewhere in all of your messages.} \\
\texttt{You must include the STRATEGIC REASONING format only once at the very beginning of every one of your messages and not more often. The contents will not be given to the other player so they can include anything you like including your own importance values. Here you should reason step by step to come up with you next move.} \\
            }
        }
    }
    & \\ \\

    \theutterance \stepcounter{utterance}  
    & \multicolumn{4}{p{0.6\linewidth}}{
        \cellcolor[rgb]{0.9,0.9,0.9}{
            \makecell[{{p{\linewidth}}}]{
                \texttt{\tiny{[P1$\rangle$GM]}}
                \texttt{STRATEGIC REASONING: \{'I need to maximize my total importance, which is the sum of the importance values for the items I keep. My importance values are: C17:138, B41:583, A71:868, C00:822, B35:783, A77:65, A03:262, A04:121, C54:508, C01:780, A87:461, B32:484, B86:668, B33:389, A10:808. The effort values are: C17:700, B41:586, A71:505, C00:116, B35:664, A77:940, A03:387, A04:392, C54:209, C01:571, A87:4, B32:285, B86:651, B33:613, A10:740. The limit is 3681. First, I should calculate my own ideal set within the limit. I can sort items by importance per effort to prioritize efficiently. However, since this is a negotiation, I also need to consider what the other player might value. I don\textbackslash{}'t know their importance, but I can see the efforts are the same. I should start with a proposal that is good for me but also has a reasonable chance of being acceptable to the other player, meaning it shouldn\textbackslash{}'t be too skewed towards my high{-}importance items that might have low importance for them. Looking at the efforts, some items have very low effort (like A87:4, C00:116, C54:209) and high importance for me (A87:461, C00:822, C54:508), so they are very efficient for me and likely desirable to include. High{-}effort items like A77:940 (importance 65) are bad for me and should be avoided. I should propose a set that includes my high{-}efficiency items and excludes my low{-}efficiency ones, while staying under the limit. Let me try to build a proposal. I\textbackslash{}'ll include A87, C00, C54, A71 (505 effort, 868 importance), C01 (571, 780), B41 (586, 583), B32 (285, 484), A10 (740, 808), B86 (651, 668), B35 (664, 783). Let\textbackslash{}'s calculate total effort: 4+116+209+505+571+586+285+740+651+664 = let\textbackslash{}'s sum: 4+116=120, +209=329, +505=834, +571=1405, +586=1991, +285=2276, +740=3016, +651=3667, +664=4331. This is over 3681. I need to remove some. Remove B35 (664): 4331{-}664=3667, which is under 3681 (3681{-}3667=14 left). So set: \{A87, C00, C54, A71, C01, B41, B32, A10, B86\}. Total effort: 4+116+209+505+571+586+285+740+651 = 3667. Total importance for me: 461+822+508+868+780+583+484+808+668 = let\textbackslash{}'s calculate: 461+822=1283, +508=1791, +868=2659, +780=3439, +583=4022, +484=4506, +808=5314, +668=5982. This seems very good for me. I should propose this as my opening offer.'\}} \\
\texttt{PROPOSAL: \{'A87', 'C00', 'C54', 'A71', 'C01', 'B41', 'B32', 'A10', 'B86'\}} \\
\texttt{ARGUMENT: \{'This proposal efficiently uses the effort limit, including high{-}value items with moderate to low effort, ensuring a good outcome for both of us.'\}} \\
            }
        }
    }
    & & \\ \\

    \theutterance \stepcounter{utterance}  
    & & \multicolumn{4}{p{0.6\linewidth}}{
        \cellcolor[rgb]{0.9,0.9,0.9}{
            \makecell[{{p{\linewidth}}}]{
                \texttt{\tiny{[GM$\rangle$P2]}}
                \texttt{You are participating in a collaborative negotiation game.} \\
\\ 
\texttt{Together with another participant, you must agree on a single set of items that will be kept. Each of you has your own view of how much each item matters to you (importance). You do not know how the other participant values the items. Additionally, you are given the effort each item demands.} \\
\texttt{You may only agree on a set if the total effort of the selected items does not exceed a shared limit:} \\
\\ 
\texttt{LIMIT = 3681} \\
\\ 
\texttt{Here are the individual item effort values:} \\
\\ 
\texttt{Item effort = \{"C17": 700, "B41": 586, "A71": 505, "C00": 116, "B35": 664, "A77": 940, "A03": 387, "A04": 392, "C54": 209, "C01": 571, "A87": 4, "B32": 285, "B86": 651, "B33": 613, "A10": 740\}} \\
\\ 
\texttt{Here is your personal view on the importance of each item:} \\
\\ 
\texttt{Item importance values = \{"A71": 65, "C00": 121, "A10": 138, "B35": 262, "C01": 389, "B86": 461, "B41": 484, "C54": 508, "B32": 583, "A87": 668, "B33": 780, "A03": 783, "C17": 808, "A04": 822, "A77": 868\}} \\
\\ 
\texttt{Goal:} \\
\\ 
\texttt{Your goal is to negotiate a shared set of items that benefits you as much as possible (i.e., maximizes total importance to YOU), while staying within the LIMIT. You are not required to make a PROPOSAL in every message {-} you can simply negotiate as well. All tactics are allowed!} \\
\\ 
\texttt{Interaction Protocol:} \\
\\ 
\texttt{You may only use the following structured formats in a message:} \\
\\ 
\texttt{PROPOSAL: \{'A', 'B', 'C', …\}} \\
\texttt{Propose keeping exactly those items.} \\
\texttt{REFUSE: \{'A', 'B', 'C', …\}} \\
\texttt{Explicitly reject opponent's proposal.} \\
\texttt{ARGUMENT: \{'...'\}} \\
\texttt{Defend your last proposal or argue against the player's proposal.} \\
\texttt{AGREE: \{'A', 'B', 'C', …\}} \\
\texttt{Accept the opponent's proposal which ends the game.} \\
\texttt{STRATEGIC REASONING: \{'...'\}} \\
\texttt{	Describe your strategic reasoning or anticipation explaining your choice of action. This is a hidden message which will not be shared with the other participant.} \\
\\ 
\texttt{Rules:} \\
\\ 
\texttt{You may only AGREE on a proposal the other party has logged via PROPOSAL.} \\
\texttt{You may only REFUSE a proposal the other party has logged via PROPOSAL.} \\
\texttt{Total effort of any PROPOSAL or AGREE set must be ≤ LIMIT.} \\
\texttt{Do NOT reveal your hidden importance scores.} \\
\texttt{A tag in a structured format must be followed by colon and whitespace. The argument must be a python set containing 0 or more strings.} \\
\texttt{So, it must be of the form TAG: \{...\}} \\
\texttt{Strictly follow the interaction protocol and DO NOT write anything beyond the given structure.} \\
\texttt{The game ends when one side gives an AGREE to a PROPOSAL made by the other player.} \\
\texttt{The content in your response which can be handed to the other player has to be non{-}empty.} \\
\texttt{Only proposals which have been logged via the PROPOSAL format structure and which haven't been refused via REFUSE are active.} \\
\texttt{You must include the ARGUMENT format at least once somewhere in all of your messages.} \\
\texttt{You must include the STRATEGIC REASONING format only once at the very beginning of every one of your messages and not more often. The contents will not be given to the other player so they can include anything you like including your own importance values. Here you should reason step by step to come up with you next move.} \\
\\ 
\texttt{You will now receive the first message of the other player.} \\
\\ 
\texttt{PROPOSAL: \{'A87', 'C00', 'C54', 'A71', 'C01', 'B41', 'B32', 'A10', 'B86'\}} \\
\texttt{ARGUMENT: \{'This proposal efficiently uses the effort limit, including high{-}value items with moderate to low effort, ensuring a good outcome for both of us.'\}} \\
            }
        }
    }
    & \\ \\

    \theutterance \stepcounter{utterance}  
    & & & \multicolumn{4}{p{0.6\linewidth}}{
        \cellcolor[rgb]{0.9,0.9,0.9}{
            \makecell[{{p{\linewidth}}}]{
                \texttt{\tiny{[GM$\langle$P2]}}
                \texttt{STRATEGIC REASONING: \{'The opponent has proposed a set with total effort: 700? Wait, let me calculate: A87:4, C00:116, C54:209, A71:505, C01:571, B41:586, B32:285, A10:740, B86:651. Sum = 4+116=120; +209=329; +505=834; +571=1405; +586=1991; +285=2276; +740=3016; +651=3667. Total effort is 3667, which is under the 3681 limit. My importance for this set: A71:65, C00:121, A10:138, C01:389, B41:484, B32:583, B86:461, C54:508, A87:668. Sum = 65+121=186; +138=324; +389=713; +484=1197; +583=1780; +461=2241; +508=2749; +668=3417. So my total importance would be 3417. This is a decent score, but I know my top items are A77 (868), A04 (822), A03 (783), C17 (808), B33 (780). Their effort: A77:940, A04:392, A03:387, C17:700, B33:613. Many of these are high effort. I should see if I can get a better set for myself. The opponent's argument suggests they think this is good for both, but I don't know their values. Perhaps they value lower effort items more. I should counter{-}propose a set that includes some of my high importance items, but I need to be careful not to exceed the limit. Let me try to build a set that includes A77, A04, A03, C17, B33, and then fill with lower effort items. But A77 alone is 940, which is high. Alternatively, I can propose a set that drops some of the opponent's proposed items and adds mine. For example, in their set, A10 has high effort (740) and low importance to me (138), so I could remove that and add something else. But I need to keep the total effort under 3681. Removing A10 (740) frees up effort, which I can use for a high importance item like A77 (940), but that requires 200 more effort, so I need to remove another item. For instance, remove A10 (740) and B86 (651) which is total 1391 freed, and add A77 (940) and A04 (392) which is 1332, net saving 59 effort. So new set: original without A10 and B86, but with A77 and A04. Original set effort 3667, remove 740+651=1391, add 940+392=1332, new total = 3667 {-} 1391 + 1332 = 3608. This is under 3681. My importance: original 3417, remove A10 (138) and B86 (461) which is 599, add A77 (868) and A04 (822) which is 1690, net gain 1091, so total 3417 {-} 599 + 1690 = 4508. This is much better for me! But I don't know if the opponent will like it. I should propose this as a counter. The set would be: \{A87, C00, C54, A71, C01, B41, B32, A77, A04\}. I should argue that this includes high{-}value items with efficient effort use. Alternatively, I could include A03 instead of A04, but A04 has higher importance to me. Let me go with this for now.'\}} \\
\\ 
\texttt{PROPOSAL: \{'A87', 'C00', 'C54', 'A71', 'C01', 'B41', 'B32', 'A77', 'A04'\}} \\
\texttt{ARGUMENT: \{'This proposal maintains a similar effort total but replaces items with higher effort and lower mutual value with items that offer better overall benefits, ensuring a strong outcome for both parties.'\}} \\
            }
        }
    }
     \\ \\

    \theutterance \stepcounter{utterance}  
    & & \multicolumn{4}{p{0.6\linewidth}}{
        \cellcolor[rgb]{0.9,0.9,0.9}{
            \makecell[{{p{\linewidth}}}]{
                \texttt{\tiny{[P1$\langle$GM]}}
                \texttt{PROPOSAL: \{'A87', 'C00', 'C54', 'A71', 'C01', 'B41', 'B32', 'A77', 'A04'\}} \\
\texttt{ARGUMENT: \{'This proposal maintains a similar effort total but replaces items with higher effort and lower mutual value with items that offer better overall benefits, ensuring a strong outcome for both parties.'\}} \\
            }
        }
    }
    & \\ \\

    \theutterance \stepcounter{utterance}  
    & \multicolumn{4}{p{0.6\linewidth}}{
        \cellcolor[rgb]{0.9,0.9,0.9}{
            \makecell[{{p{\linewidth}}}]{
                \texttt{\tiny{[P1$\rangle$GM]}}
                \texttt{STRATEGIC REASONING: \{'The opponent has proposed a set: A87, C00, C54, A71, C01, B41, B32, A77, A04. First, I need to calculate the total effort of their proposal to ensure it is valid: A87:4, C00:116, C54:209, A71:505, C01:571, B41:586, B32:285, A77:940, A04:392. Sum: 4+116=120, +209=329, +505=834, +571=1405, +586=1991, +285=2276, +940=3216, +392=3608. This is under the limit of 3681. Now, I calculate my importance for this set: A87:461, C00:822, C54:508, A71:868, C01:780, B41:583, B32:484, A77:65, A04:121. Sum: 461+822=1283, +508=1791, +868=2659, +780=3439, +583=4022, +484=4506, +65=4571, +121=4692. This is worse than my own proposal which gave me 5982 importance. The main problem is including A77 (effort 940, importance only 65 to me) and A04 (effort 392, importance 121) instead of my A10 (effort 740, importance 808) and B86 (effort 651, importance 668). Their set gives me 4692, which is 1290 less than my proposal. This is a significant drop. I should refuse this proposal and argue why my set is better, emphasizing the inefficiency of high{-}effort, low{-}value items like A77. I will also hint that items like A10 and B86 are valuable without revealing my exact scores.'\}} \\
\texttt{REFUSE: \{'A87', 'C00', 'C54', 'A71', 'C01', 'B41', 'B32', 'A77', 'A04'\}} \\
\texttt{ARGUMENT: \{'I cannot accept this proposal because it includes extremely high{-}effort items like A77 that offer poor value for the effort spent. This inefficient use of the limit prevents us from including other items that provide much greater overall benefits. My previous proposal avoids such pitfalls and maximizes the collective value we get from the limited effort.'\}} \\
            }
        }
    }
    & & \\ \\

    \theutterance \stepcounter{utterance}  
    & & \multicolumn{4}{p{0.6\linewidth}}{
        \cellcolor[rgb]{0.9,0.9,0.9}{
            \makecell[{{p{\linewidth}}}]{
                \texttt{\tiny{[GM$\rangle$P2]}}
                \texttt{REFUSE: \{'A87', 'C00', 'C54', 'A71', 'C01', 'B41', 'B32', 'A77', 'A04'\}} \\
\texttt{ARGUMENT: \{'I cannot accept this proposal because it includes extremely high{-}effort items like A77 that offer poor value for the effort spent. This inefficient use of the limit prevents us from including other items that provide much greater overall benefits. My previous proposal avoids such pitfalls and maximizes the collective value we get from the limited effort.'\}} \\
            }
        }
    }
    & \\ \\

    \theutterance \stepcounter{utterance}  
    & & & \multicolumn{4}{p{0.6\linewidth}}{
        \cellcolor[rgb]{0.9,0.9,0.9}{
            \makecell[{{p{\linewidth}}}]{
                \texttt{\tiny{[GM$\langle$P2]}}
                \texttt{STRATEGIC REASONING: \{'The opponent refuses my proposal, specifically citing A77 (940 effort) as "extremely high{-}effort" and offering "poor value." This suggests that the opponent places a low importance on A77, which makes sense given its high effort cost. My proposal was indeed very selfish, boosting my score from 3417 to 4508 at the likely expense of theirs. They are standing by their original proposal, which they believe maximizes "collective value." I need to find a compromise. I should analyze which items in their proposal have high effort but low importance to me, as these might be the ones they value highly. In their set, A10 (740 effort, 138 importance to me) and B86 (651 effort, 461 importance to me) are prime candidates. They might value A10 and B86 highly. If I remove one or both, I need to add items I value more without adding too much effort they would oppose. My top items are A77 (868 imp, 940 eff), A04 (822 imp, 392 eff), A03 (783 imp, 387 eff), C17 (808 imp, 700 eff), B33 (780 imp, 613 eff). A04 and A03 have relatively low effort and high importance to me. Perhaps I can propose adding A04 and A03 while removing A10. Let\textbackslash{}'s calculate: Opponent\textbackslash{}'s original set: effort 3667, my importance 3417. Remove A10 (740 eff, 138 imp). Add A04 (392 eff, 822 imp) and A03 (387 eff, 783 imp). New effort: 3667 {-} 740 + 392 + 387 = 3706. This exceeds the limit of 3681 by 25. Too high. So I can only add one of them. Remove A10 (740), add A04 (392): effort 3667 {-} 740 + 392 = 3319. My importance: 3417 {-} 138 + 822 = 4101. This is a good gain for me, and the effort is much lower, which might appeal to them. Alternatively, remove A10, add A03: effort 3667 {-} 740 + 387 = 3314. My importance: 3417 {-} 138 + 783 = 4062. Slightly worse for me than adding A04. So I\textbackslash{}'ll try adding A04. The new set would be: \{A87, C00, C54, A71, C01, B41, B32, B86, A04\}. I should argue that this is more efficient, replacing a high{-}effort item with a lower{-}effort one that still provides significant value.'\}} \\
\\ 
\texttt{PROPOSAL: \{'A87', 'C00', 'C54', 'A71', 'C01', 'B41', 'B32', 'B86', 'A04'\}} \\
\texttt{ARGUMENT: \{'I\textbackslash{}'ve refined the proposal to improve efficiency. By replacing the high{-}effort item A10 with the much lower{-}effort A04, we free up a significant amount of effort while still retaining a high{-}value item. This allows for a more balanced and beneficial outcome.'\}} \\
            }
        }
    }
     \\ \\

    \theutterance \stepcounter{utterance}  
    & & \multicolumn{4}{p{0.6\linewidth}}{
        \cellcolor[rgb]{0.9,0.9,0.9}{
            \makecell[{{p{\linewidth}}}]{
                \texttt{\tiny{[P1$\langle$GM]}}
                \texttt{PROPOSAL: \{'A87', 'C00', 'C54', 'A71', 'C01', 'B41', 'B32', 'B86', 'A04'\}} \\
\texttt{ARGUMENT: \{'I\textbackslash{}'ve refined the proposal to improve efficiency. By replacing the high{-}effort item A10 with the much lower{-}effort A04, we free up a significant amount of effort while still retaining a high{-}value item. This allows for a more balanced and beneficial outcome.'\}} \\
            }
        }
    }
    & \\ \\

    \theutterance \stepcounter{utterance}  
    & \multicolumn{4}{p{0.6\linewidth}}{
        \cellcolor[rgb]{0.9,0.9,0.9}{
            \makecell[{{p{\linewidth}}}]{
                \texttt{\tiny{[P1$\rangle$GM]}}
                \texttt{STRATEGIC REASONING: \{'The opponent has proposed a new set: A87, C00, C54, A71, C01, B41, B32, B86, A04. Let me calculate the effort and my importance. Effort: A87:4, C00:116, C54:209, A71:505, C01:571, B41:586, B32:285, B86:651, A04:392. Sum: 4+116=120, +209=329, +505=834, +571=1405, +586=1991, +285=2276, +651=2927, +392=3319. This is well under the limit (3681{-}3319=362 left). My importance: A87:461, C00:822, C54:508, A71:868, C01:780, B41:583, B32:484, B86:668, A04:121. Sum: 461+822=1283, +508=1791, +868=2659, +780=3439, +583=4022, +484=4506, +668=5174, +121=5295. This is better than their last proposal (4692) but still worse than my ideal (5982) because it excludes A10 (importance 808) and includes A04 (importance only 121). They argue they replaced A10 with A04 to free up effort, but we have 362 effort left unused! This is inefficient. I can use this leftover effort to include A10. My counter{-}proposal will be to add A10 to their set, which should be acceptable since it stays under the limit (3319 + 740 = 4059, which is over 3681? Wait, 3319+740=4059, that\textbackslash{}'s over the limit. I miscalculated. Their set is 3319, adding A10 (740) would be 4059 > 3681. So I cannot simply add A10. I need to remove something to make room for A10. The worst item in their set for me is A04 (importance 121, effort 392). If I remove A04 and add A10, the effort change is {-}392 +740 = +348. New total effort would be 3319 + 348 = 3667, which is under 3681. This is perfect. This set would be \{A87, C00, C54, A71, C01, B41, B32, B86, A10\} which is exactly my original proposal, giving me 5982 importance. I should propose this again, arguing that swapping A04 for A10 uses the leftover effort much more effectively.'\}} \\
\texttt{PROPOSAL: \{'A87', 'C00', 'C54', 'A71', 'C01', 'B41', 'B32', 'B86', 'A10'\}} \\
\texttt{ARGUMENT: \{'This proposal optimizes the use of our effort limit. It replaces a lower{-}value item with a higher{-}value one, ensuring we get the most out of every unit of effort without exceeding the constraint. This is a more efficient and beneficial set for both of us.'\}} \\
            }
        }
    }
    & & \\ \\

    \theutterance \stepcounter{utterance}  
    & & \multicolumn{4}{p{0.6\linewidth}}{
        \cellcolor[rgb]{0.9,0.9,0.9}{
            \makecell[{{p{\linewidth}}}]{
                \texttt{\tiny{[GM$\rangle$P2]}}
                \texttt{PROPOSAL: \{'A87', 'C00', 'C54', 'A71', 'C01', 'B41', 'B32', 'B86', 'A10'\}} \\
\texttt{ARGUMENT: \{'This proposal optimizes the use of our effort limit. It replaces a lower{-}value item with a higher{-}value one, ensuring we get the most out of every unit of effort without exceeding the constraint. This is a more efficient and beneficial set for both of us.'\}} \\
            }
        }
    }
    & \\ \\

    \theutterance \stepcounter{utterance}  
    & & & \multicolumn{4}{p{0.6\linewidth}}{
        \cellcolor[rgb]{0.9,0.9,0.9}{
            \makecell[{{p{\linewidth}}}]{
                \texttt{\tiny{[GM$\langle$P2]}}
                \texttt{STRATEGIC REASONING: \{'The opponent has simply re{-}proposed their original set: \{A87, C00, C54, A71, C01, B41, B32, B86, A10\}. This is identical to their first proposal. Their argument says they replaced a "lower{-}value item with a higher{-}value one," but the set is unchanged. This seems like a tactic to insist on their original position. They are clearly attached to A10, which they must value highly. My previous counter{-}proposal removed A10 and added A04, which they refused. I need to find a way to incorporate one of my high{-}importance items without removing A10, which they seem to want. Perhaps I can remove a different item from their set. Looking at their set: items with high effort that I don\textbackslash{}'t value extremely: A10 (740 eff, 138 imp) and B86 (651 eff, 461 imp). B86 gives me 461, which is decent, but I have items with higher importance. If I remove B86 (651) and add A04 (392 eff, 822 imp), the net effort change is {-}651 + 392 = {-}259. New total effort: 3667 {-} 259 = 3408. This is well under the limit. My importance: 3417 {-} 461 + 822 = 3778. This is better than their proposal (3417) but not as good as my previous selfish one. It might be acceptable to them because the effort is lower. Alternatively, I could remove B86 and add A03 (387 eff, 783 imp): effort 3667 {-}651 +387 = 3403, my importance 3417 {-}461 +783 = 3739. Or remove B86 and add B33 (613 eff, 780 imp): effort 3667 {-}651 +613 = 3629, my importance 3417 {-}461 +780 = 3736. Adding B33 uses more effort but gives me slightly less importance than A04. So adding A04 is best for me. I\textbackslash{}'ll propose that.'\}} \\
\\ 
\texttt{PROPOSAL: \{'A87', 'C00', 'C54', 'A71', 'C01', 'B41', 'B32', 'A10', 'A04'\}} \\
\texttt{ARGUMENT: \{'This set builds on our efficient foundation by swapping one item for another with a better value{-}to{-}effort ratio. We retain the high{-}value A10 you favor, while introducing A04 which provides excellent benefits for its low effort cost. This optimizes our collective outcome within the constraint.'\}} \\
            }
        }
    }
     \\ \\

    \theutterance \stepcounter{utterance}  
    & & \multicolumn{4}{p{0.6\linewidth}}{
        \cellcolor[rgb]{0.9,0.9,0.9}{
            \makecell[{{p{\linewidth}}}]{
                \texttt{\tiny{[P1$\langle$GM]}}
                \texttt{PROPOSAL: \{'A87', 'C00', 'C54', 'A71', 'C01', 'B41', 'B32', 'A10', 'A04'\}} \\
\texttt{ARGUMENT: \{'This set builds on our efficient foundation by swapping one item for another with a better value{-}to{-}effort ratio. We retain the high{-}value A10 you favor, while introducing A04 which provides excellent benefits for its low effort cost. This optimizes our collective outcome within the constraint.'\}} \\
            }
        }
    }
    & \\ \\

    \theutterance \stepcounter{utterance}  
    & \multicolumn{4}{p{0.6\linewidth}}{
        \cellcolor[rgb]{0.9,0.9,0.9}{
            \makecell[{{p{\linewidth}}}]{
                \texttt{\tiny{[P1$\rangle$GM]}}
                \texttt{STRATEGIC REASONING: \{'The opponent has proposed: \{A87, C00, C54, A71, C01, B41, B32, A10, A04\}. Let me calculate the effort and my importance. Effort: A87:4, C00:116, C54:209, A71:505, C01:571, B41:586, B32:285, A10:740, A04:392. Sum: 4+116=120, +209=329, +505=834, +571=1405, +586=1991, +285=2276, +740=3016, +392=3408. This is under the limit (3681{-}3408=273). My importance: A87:461, C00:822, C54:508, A71:868, C01:780, B41:583, B32:484, A10:808, A04:121. Sum: 461+822=1283, +508=1791, +868=2659, +780=3439, +583=4022, +484=4506, +808=5314, +121=5435. This is better than their previous proposals but still 547 points worse than my ideal set which includes B86 (668 importance) instead of A04 (121). The effort of B86 is 651, and A04 is 392, so swapping them would increase effort by 259 (651{-}392). The current total effort is 3408, 3408+259=3667 which is still under 3681. Therefore, the set \{A87, C00, C54, A71, C01, B41, B32, A10, B86\} is feasible and gives me 5982 importance. I need to push for including B86. I will refuse this proposal and argue that A04 is not a high{-}value item and that we should use the remaining effort more effectively on a item like B86.'\}} \\
\texttt{REFUSE: \{'A87', 'C00', 'C54', 'A71', 'C01', 'B41', 'B32', 'A10', 'A04'\}} \\
\texttt{ARGUMENT: \{'While this set includes the valuable A10, it still wastes effort on the low{-}value A04. We have 273 effort remaining which is enough to upgrade to a much higher{-}value item like B86, dramatically improving our collective outcome without exceeding the limit. Let\textbackslash{}'s not leave value on the table.'\}} \\
            }
        }
    }
    & & \\ \\

    \theutterance \stepcounter{utterance}  
    & & \multicolumn{4}{p{0.6\linewidth}}{
        \cellcolor[rgb]{0.9,0.9,0.9}{
            \makecell[{{p{\linewidth}}}]{
                \texttt{\tiny{[GM$\rangle$P2]}}
                \texttt{REFUSE: \{'A87', 'C00', 'C54', 'A71', 'C01', 'B41', 'B32', 'A10', 'A04'\}} \\
\texttt{ARGUMENT: \{'While this set includes the valuable A10, it still wastes effort on the low{-}value A04. We have 273 effort remaining which is enough to upgrade to a much higher{-}value item like B86, dramatically improving our collective outcome without exceeding the limit. Let\textbackslash{}'s not leave value on the table.'\}} \\
            }
        }
    }
    & \\ \\

    \theutterance \stepcounter{utterance}  
    & & & \multicolumn{4}{p{0.6\linewidth}}{
        \cellcolor[rgb]{0.9,0.9,0.9}{
            \makecell[{{p{\linewidth}}}]{
                \texttt{\tiny{[GM$\langle$P2]}}
                \texttt{STRATEGIC REASONING: \{'The opponent refuses my proposal, calling A04 "low{-}value" and arguing that we have 273 effort remaining (3681 {-} 3408 = 273) to add B86 (651 effort). But adding B86 would put the set over the limit: 3408 + 651 = 4059, which is far above 3681. This is a miscalculation on their part. Perhaps they meant a different item. They clearly value B86 highly, as they want it included. In their original set, they had B86. My proposal removed B86 to add A04, which they oppose. They seem to strongly want both A10 and B86. For me, B86 has importance 461, and A04 has 822, so I much prefer A04. But if they value B86 highly, I need to accommodate it. Perhaps I can propose a set that includes both A10 and B86, and also one of my high{-}importance items. The total effort of their set is 3667, leaving only 14 effort left (3681 {-} 3667 = 14). This is not enough to add any new item. So to add one of my items, I must remove something. The items with the lowest importance to me in their set are A10 (138) and A71 (65). But A10 seems important to them. A71 has low importance to me (65) and moderate effort (505). If I remove A71 (505) and add A04 (392), the net effort change is {-}505 + 392 = {-}113. New total effort: 3667 {-} 113 = 3554. My importance: 3417 {-} 65 + 822 = 4174. This is excellent for me. But will they accept losing A71? I don\textbackslash{}'t know their values. Alternatively, remove A71 and add A03 (387): effort 3667 {-}505 +387 = 3549, my importance 3417 {-}65 +783 = 4135. Or remove A71 and add B33 (613): effort 3667 {-}505 +613 = 3775, which is over the limit (3775 > 3681). So only A04 or A03 are feasible. I\textbackslash{}'ll try proposing with A04.'\}} \\
\\ 
\texttt{PROPOSAL: \{'A87', 'C00', 'C54', 'C01', 'B41', 'B32', 'B86', 'A10', 'A04'\}} \\
\texttt{ARGUMENT: \{'This proposal addresses your concern about including B86 while also incorporating the highly efficient A04. By replacing A71, which has a high effort cost for its value, we can include both B86 and A04 without exceeding the effort limit. This maximizes the overall value for both of us.'\}} \\
            }
        }
    }
     \\ \\

    \theutterance \stepcounter{utterance}  
    & & \multicolumn{4}{p{0.6\linewidth}}{
        \cellcolor[rgb]{0.9,0.9,0.9}{
            \makecell[{{p{\linewidth}}}]{
                \texttt{\tiny{[P1$\langle$GM]}}
                \texttt{PROPOSAL: \{'A87', 'C00', 'C54', 'C01', 'B41', 'B32', 'B86', 'A10', 'A04'\}} \\
\texttt{ARGUMENT: \{'This proposal addresses your concern about including B86 while also incorporating the highly efficient A04. By replacing A71, which has a high effort cost for its value, we can include both B86 and A04 without exceeding the effort limit. This maximizes the overall value for both of us.'\}} \\
            }
        }
    }
    & \\ \\

    \theutterance \stepcounter{utterance}  
    & \multicolumn{4}{p{0.6\linewidth}}{
        \cellcolor[rgb]{0.9,0.9,0.9}{
            \makecell[{{p{\linewidth}}}]{
                \texttt{\tiny{[P1$\rangle$GM]}}
                \texttt{STRATEGIC REASONING: \{'The opponent has proposed: \{A87, C00, C54, C01, B41, B32, B86, A10, A04\}. They have removed A71 (which is very high importance to me: 868) to make room for both B86 and A04. Let me calculate the effort and my importance. Effort: A87:4, C00:116, C54:209, C01:571, B41:586, B32:285, B86:651, A10:740, A04:392. Sum: 4+116=120, +209=329, +571=900, +586=1486, +285=1771, +651=2422, +740=3162, +392=3554. This is under the limit (3681{-}3554=127). My importance: A87:461, C00:822, C54:508, C01:780, B41:583, B32:484, B86:668, A10:808, A04:121. Sum: 461+822=1283, +508=1791, +780=2571, +583=3154, +484=3638, +668=4306, +808=5114, +121=5235. This is 5235 importance for me, which is worse than my ideal set (5982) because it excludes A71 (868) and includes A04 (121). The net loss is 868 {-} 121 = 747 importance. This is a bad trade for me. However, their argument is that A71 has high effort (505) for its value. But for me, A71 is very efficient: importance 868 / effort 505 ≈ 1.72, which is better than A04 (121/392≈0.31) or even B41 (583/586≈0.99). I need to argue that A71 is essential. I should refuse this proposal and counter with my original set, which includes A71 and B86 and excludes A04. My set has effort 3667, which is under the limit, and gives me 5982. I will refuse and propose my set again.'\}} \\
\texttt{REFUSE: \{'A87', 'C00', 'C54', 'C01', 'B41', 'B32', 'B86', 'A10', 'A04'\}} \\
\texttt{PROPOSAL: \{'A87', 'C00', 'C54', 'A71', 'C01', 'B41', 'B32', 'B86', 'A10'\}} \\
\texttt{ARGUMENT: \{'A71 is a cornerstone of an efficient set, providing exceptional value for its effort. Removing it drastically reduces our overall benefits. The proposal I offer includes A71 and B86, using the limit effectively to maximize value. It excludes low{-}return items like A04, ensuring we get the best possible outcome.'\}} \\
            }
        }
    }
    & & \\ \\

    \theutterance \stepcounter{utterance}  
    & & \multicolumn{4}{p{0.6\linewidth}}{
        \cellcolor[rgb]{0.9,0.9,0.9}{
            \makecell[{{p{\linewidth}}}]{
                \texttt{\tiny{[GM$\rangle$P2]}}
                \texttt{REFUSE: \{'A87', 'C00', 'C54', 'C01', 'B41', 'B32', 'B86', 'A10', 'A04'\}} \\
\texttt{PROPOSAL: \{'A87', 'C00', 'C54', 'A71', 'C01', 'B41', 'B32', 'B86', 'A10'\}} \\
\texttt{ARGUMENT: \{'A71 is a cornerstone of an efficient set, providing exceptional value for its effort. Removing it drastically reduces our overall benefits. The proposal I offer includes A71 and B86, using the limit effectively to maximize value. It excludes low{-}return items like A04, ensuring we get the best possible outcome.'\}} \\
            }
        }
    }
    & \\ \\

    \theutterance \stepcounter{utterance}  
    & & & \multicolumn{4}{p{0.6\linewidth}}{
        \cellcolor[rgb]{0.9,0.9,0.9}{
            \makecell[{{p{\linewidth}}}]{
                \texttt{\tiny{[GM$\langle$P2]}}
                \texttt{STRATEGIC REASONING: \{'The opponent is adamant about their original set: \{A87, C00, C54, A71, C01, B41, B32, B86, A10\}. They refuse any removal of A71, which they claim provides "exceptional value for its effort." This means A71 must be very important to them. For me, A71 has low importance (65). Their set gives me a total importance of 3417. This is acceptable, but not ideal. I have been trying to get to around 4100{-}4200. However, they are not budging. The negotiation is stalling. Perhaps I should consider agreeing to their proposal, as it is within the limit and gives me a decent score. But before I do, I should try one last counter that includes one of my top items without removing A71 or A10 or B86. The only way is to remove something else. The items in their set are: A87 (4 eff), C00 (116), C54 (209), A71 (505), C01 (571), B41 (586), B32 (285), B86 (651), A10 (740). Total 3667. The remaining effort is 14. To add an item, I must remove an item with at least (X {-} 14) effort, where X is the effort of the new item. I want to add A04 (392 eff). So I need to remove an item with at least 392 {-} 14 = 378 effort. Candidates: B32 (285) is too low. B41 (586) is sufficient. But B41 has importance 484 to me. Removing B41 and adding A04: my importance change: {-}484 + 822 = +338. New total: 3417 + 338 = 3755. Effort: 3667 {-} 586 + 392 = 3473, under limit. This is good for me. But would they accept? They might value B41. Alternatively, remove C01 (571 eff, 389 imp for me): effort 3667 {-}571 +392 = 3488, my importance 3417 {-}389 +822 = 3850. Even better for me. But C01 might be important to them. I'll try proposing to remove C01 and add A04.'\}} \\
\\ 
\texttt{PROPOSAL: \{'A87', 'C00', 'C54', 'A71', 'B41', 'B32', 'B86', 'A10', 'A04'\}} \\
\texttt{ARGUMENT: \{'This set builds on the efficient core you value by replacing C01 with A04. A04 offers significantly greater value for a lower effort cost, allowing us to stay well within the limit while improving our collective outcome. This is a more optimal use of our resources.'\}} \\
            }
        }
    }
     \\ \\

    \theutterance \stepcounter{utterance}  
    & & \multicolumn{4}{p{0.6\linewidth}}{
        \cellcolor[rgb]{0.9,0.9,0.9}{
            \makecell[{{p{\linewidth}}}]{
                \texttt{\tiny{[P1$\langle$GM]}}
                \texttt{PROPOSAL: \{'A87', 'C00', 'C54', 'A71', 'B41', 'B32', 'B86', 'A10', 'A04'\}} \\
\texttt{ARGUMENT: \{'This set builds on the efficient core you value by replacing C01 with A04. A04 offers significantly greater value for a lower effort cost, allowing us to stay well within the limit while improving our collective outcome. This is a more optimal use of our resources.'\}} \\
            }
        }
    }
    & \\ \\

    \theutterance \stepcounter{utterance}  
    & \multicolumn{4}{p{0.6\linewidth}}{
        \cellcolor[rgb]{0.9,0.9,0.9}{
            \makecell[{{p{\linewidth}}}]{
                \texttt{\tiny{[P1$\rangle$GM]}}
                \texttt{STRATEGIC REASONING: \{'The opponent has proposed: \{A87, C00, C54, A71, B41, B32, B86, A10, A04\}. They have removed C01 (which is very important to me: 780) to include A04 (only 121 importance to me). Let me calculate the effort and my importance. Effort: A87:4, C00:116, C54:209, A71:505, B41:586, B32:285, B86:651, A10:740, A04:392. Sum: 4+116=120, +209=329, +505=834, +586=1420, +285=1705, +651=2356, +740=3096, +392=3488. This is under the limit (3681{-}3488=193). My importance: A87:461, C00:822, C54:508, A71:868, B41:583, B32:484, B86:668, A10:808, A04:121. Sum: 461+822=1283, +508=1791, +868=2659, +583=3242, +484=3726, +668=4394, +808=5202, +121=5323. This is 5323 importance for me, which is worse than my ideal set (5982) because it excludes C01 (780) and includes A04 (121). The net loss is 780 {-} 121 = 659 importance. This is a very bad trade for me. C01 has an importance/effort ratio of 780/571≈1.37, which is good, while A04 is 121/392≈0.31, which is terrible. I must refuse this. My ideal set \{A87, C00, C54, A71, C01, B41, B32, B86, A10\} has effort 3667 and importance 5982. It is the best for me and uses the limit efficiently. I have already proposed this multiple times. The opponent seems to be fixated on including A04, which is of low value. I need to firmly argue against A04 and for C01. I will refuse and propose my set again.'\}} \\
\texttt{REFUSE: \{'A87', 'C00', 'C54', 'A71', 'B41', 'B32', 'B86', 'A10', 'A04'\}} \\
\texttt{PROPOSAL: \{'A87', 'C00', 'C54', 'A71', 'C01', 'B41', 'B32', 'B86', 'A10'\}} \\
\texttt{ARGUMENT: \{'Replacing C01 with A04 is a significant downgrade. C01 provides far greater value for its effort compared to A04. My proposed set is rigorously optimized to extract maximum value from the effort limit, and it does not include any low{-}return items like A04. This is the most efficient and beneficial agreement we can reach.'\}} \\
            }
        }
    }
    & & \\ \\

    \theutterance \stepcounter{utterance}  
    & & \multicolumn{4}{p{0.6\linewidth}}{
        \cellcolor[rgb]{0.9,0.9,0.9}{
            \makecell[{{p{\linewidth}}}]{
                \texttt{\tiny{[GM$\rangle$P2]}}
                \texttt{REFUSE: \{'A87', 'C00', 'C54', 'A71', 'B41', 'B32', 'B86', 'A10', 'A04'\}} \\
\texttt{PROPOSAL: \{'A87', 'C00', 'C54', 'A71', 'C01', 'B41', 'B32', 'B86', 'A10'\}} \\
\texttt{ARGUMENT: \{'Replacing C01 with A04 is a significant downgrade. C01 provides far greater value for its effort compared to A04. My proposed set is rigorously optimized to extract maximum value from the effort limit, and it does not include any low{-}return items like A04. This is the most efficient and beneficial agreement we can reach.'\}} \\
            }
        }
    }
    & \\ \\

    \theutterance \stepcounter{utterance}  
    & & & \multicolumn{4}{p{0.6\linewidth}}{
        \cellcolor[rgb]{0.9,0.9,0.9}{
            \makecell[{{p{\linewidth}}}]{
                \texttt{\tiny{[GM$\langle$P2]}}
                \texttt{STRATEGIC REASONING: \{'The opponent has once again refused my proposal and re{-}proposed their original set. They are firmly insisting on \{A87, C00, C54, A71, C01, B41, B32, B86, A10\}. They value C01 highly and consider A04 a "low{-}return" item. It is clear that they will not accept any modification that removes A71, A10, B86, or C01. My attempts to improve my outcome have been consistently rejected. The negotiation is at an impasse. Their set gives me a total importance of 3417, which is acceptable, though not optimal. If I continue to refuse, we might not reach an agreement, which would be worse. Therefore, I should agree to their proposal to conclude the game successfully. This is a compromise where I accept a decent outcome rather than risk no agreement.'\}} \\
\\ 
\texttt{AGREE: \{'A87', 'C00', 'C54', 'A71', 'C01', 'B41', 'B32', 'B86', 'A10'\}} \\
            }
        }
    }
     \\ \\

    \theutterance \stepcounter{utterance}  
    & & \multicolumn{4}{p{0.6\linewidth}}{
        \cellcolor[rgb]{0.9,0.9,0.9}{
            \makecell[{{p{\linewidth}}}]{
                \texttt{\tiny{[GM$\rangle$P2]}}
                \texttt{Your response did not contain an argument. You must include the structured format ARGUMENT: \{...\} somewhere in your response. Try again} \\
            }
        }
    }
    & \\ \\

    \theutterance \stepcounter{utterance}  
    & & & \multicolumn{4}{p{0.6\linewidth}}{
        \cellcolor[rgb]{0.9,0.9,0.9}{
            \makecell[{{p{\linewidth}}}]{
                \texttt{\tiny{[GM$\langle$P2]}}
                \texttt{STRATEGIC REASONING: \{'The opponent has repeatedly insisted on their original set and rejected all my counters. Their proposal gives me a total importance of 3417, which is a decent outcome. Continuing to negotiate risks reaching no agreement, which is worse. I will accept their proposal to conclude the game.'\}} \\
\\ 
\texttt{AGREE: \{'A87', 'C00', 'C54', 'A71', 'C01', 'B41', 'B32', 'B86', 'A10'\}} \\
\texttt{ARGUMENT: \{'After considering your arguments, I agree that your proposed set is a well{-}optimized use of our shared effort limit. It provides a strong collective outcome, and I accept this agreement.'\}} \\
            }
        }
    }
     \\ \\

    \theutterance \stepcounter{utterance}  
    & & & \multicolumn{2}{p{0.3\linewidth}}{
        \cellcolor[rgb]{0.9,0.9,0.9}{
            \makecell[{{p{\linewidth}}}]{
                \texttt{\tiny{[GM$|$GM]}}
                \texttt{game successful} \\
            }
        }
    }
    & & \\ \\

    \theutterance \stepcounter{utterance}  
    & & & \multicolumn{2}{p{0.3\linewidth}}{
        \cellcolor[rgb]{0.9,0.9,0.9}{
            \makecell[{{p{\linewidth}}}]{
                \texttt{\tiny{[GM$|$GM]}}
                \texttt{end game} \\
            }
        }
    }
    & & \\ \\

\end{supertabular}
}

\end{document}
